% frontmatter/firstwords/como-usar.tex -- instrucciones para el adulto,
% en inglés como todo el cuaderno (ver notes/06-first-words.md). Es la
% página del cuaderno de primeras palabras en español
% (frontmatter/palabras/como-usar.tex), con lo que cambia al leer en
% inglés: se lee por sonidos, no por sílabas, y cada sonido tiene su
% botón.
\section*{How to use this book}

This book is for a child who is starting to read in English: who knows
the letters, or nearly all of them, and is beginning to put their sounds
together. There is one page for every weekday of the year, holidays
too; five or ten minutes a day are enough. Every page has the same
parts:

\begin{itemize}[leftmargin=6mm]
\item \textbf{Date}, \textbf{Read} and \textbf{Notes}, for the grown-up:
  tick every day how the reading went. After a year, it shows how far
  the child has come.
\item \textbf{The word of the day}, in the green box, with a
  \textbf{sound button} under each sound: a dot for a sound written with
  one letter (\textit{c, a, t}), a line for two or three letters
  (\textit{sh, ee, igh}), and an arc when two letters far apart make one
  sound together (the magic e of \textit{cake}). The child touches each
  button and says its sound, then says the sounds faster and faster
  until the word comes: \textit{c -- a -- t\ldots\ cat!} One word a day
  in autumn, two in winter, three in spring, a sentence in summer.
\item \textbf{An activity}, which nearly always ends with a picture to
  colour: the picture of a word the child has just read. The
  instructions (small and grey) are \textbf{for the grown-up to read
  aloud}: the child only reads what is printed big.
\end{itemize}

\subsection*{Sounds, not letter names}

{\small
\begin{itemize}[leftmargin=6mm, itemsep=1pt, topsep=2pt]
\item Say the sounds, not the names of the letters (\textit{sss}, not
  \textit{ess}), and keep them short: \textit{mmm}, not \textit{muh}, or
  \textit{c-a-t} turns into \textit{cuh-a-tuh}.
\item English vowels are not Spanish vowels: \textit{a} as in
  \textit{cat}, \textit{e} as in \textit{hen}, \textit{i} as in
  \textit{pin}, \textit{o} as in \textit{dog}, \textit{u} as in
  \textit{sun}. The next page has every sound in the book, each with a
  word: a new sound is worth a minute before the word of the day.
\item The names that can't be read sound by sound yet (\textit{Lucía,
  Toby, Amy}\ldots) have no buttons: they are read whole, like the
  child's own name. So are the \textbf{tricky words} (\textit{the, said,
  you}\ldots), which don't sound the way they are written: they come a
  few at a time, in a frame, on the spring Fridays, and the summer
  sentences use them.
\end{itemize}}

\subsection*{The activities}

{\small
\begin{itemize}[leftmargin=6mm, itemsep=1pt, topsep=2pt]
\item \textbf{\lblColorea} (or \textbf{\lblDibuja}): the simplest way to
  check that a word has been understood, not only sounded out.
\item \textbf{\lblCompleta}: colour it and add what is missing (a tail,
  a star), or turn a few lines into a drawing.
\item \textbf{\lblTraza}: trace a letter of the word, then write it (all
  26 come during the year). \textbf{\lblEscribe}: the same, with the
  whole word.
\item \textbf{\lblEncuentra}: circle the word among others that look
  very like it (\textit{pen, pin, pan}\ldots): every letter counts.
\item \textbf{\lblPalmadas}: read each word one sound at a time and
  colour a circle for each sound (the answer key has the numbers).
\item \textbf{\lblAdivina}: the grown-up reads a riddle; the child
  colours or draws the answer (from winter, and writes it).
\item \textbf{\lblUne}: join each word to its picture or to the same word
  in CAPITALS (in summer, other pairs: opposites, rhymes\ldots).
\item \textbf{\lblSiNo} (summer only): decide whether very short
  sentences make sense.
\item \textbf{\lblRelee} and \textbf{\lblRepasa} (Fridays): the words of
  the week, to read again and tick the ones that come on their own.
\end{itemize}}

\subsection*{How to help}

{\small
\begin{itemize}[leftmargin=6mm, itemsep=1pt, topsep=2pt]
\item Stuck on a word? Point to the buttons one by one and say the
  sounds together, then let the child say the word. If it doesn't come,
  say it yourself and let them repeat it -- better that than a fight.
\item No need to correct every mistake: if they read \textit{pin} for
  \textit{pen}, ask ``Are you sure? Look at the middle.''
\item If a whole week goes well at the first try, two pages a day is
  fine; if it is hard, read the page of the day before again. And the
  activity can wait: the reading is what matters.
\end{itemize}}
\newpage
